% META: {"id": "1SPE-TRIGO-EV-B-corrige", "chapitre": "1SPE-TRIGONOMETRIE", "type_objet": "corrige_evaluation", "evaluation_ref": "1SPE-TRIGO-EV-B", "version": "B", "capacites_codes": ["C1","C2"], "sources_inspiration": [], "parametres_sympy": {}, "fichier_tex": "chapitres/1SPE-TRIGONOMETRIE/evaluations/1SPE-TRIGO-EV-B-corrige.tex", "status": "approved"}
% BEGIN-VERIFY
% from sympy import *
% x = symbols('x', real=True)
% assert 315 * pi / 180 == 7*pi/4
% assert 5*pi/3 * 180/pi == 300
% assert Rational(-13,6) + 2 == Rational(-1,6)
% # But -1/6 is not in ]-pi;pi] as fraction of pi... let me recalculate
% # -13pi/6 + 2pi = -13pi/6 + 12pi/6 = -pi/6
% assert cos(3*pi/4) == -sqrt(2)/2
% assert sin(3*pi/4) == sqrt(2)/2
% assert cos(4*pi/3) == Rational(-1,2)
% assert sin(4*pi/3) == -sqrt(3)/2
% assert cos(11*pi/6) == sqrt(3)/2
% assert Rational(5,13)**2 + Rational(12,13)**2 == 1
% assert 8 * 3*pi/4 == 6*pi
% assert Rational(7, 4) == Rational(7, 4)
% assert (5*pi/3) / (2*pi) == Rational(5, 6)
% assert cos(4*pi/3) == Rational(-1,2) and sin(4*pi/3) == -sqrt(3)/2
% assert cos(7*pi/4) == sqrt(2)/2 and sin(7*pi/4) == -sqrt(2)/2
% assert 20 + 18*sin(3*pi/4) == 20 + 9*sqrt(2)
% assert abs(float(20 + 9*sqrt(2)) - 32.7) < 0.05
% END-VERIFY

\begin{corrige}{1SPE-TRIGO-EV-B}

%===============================================================
\section*{Exercice 1 \hfill (4 points) — C1}
%===============================================================

\baremeIndicatif{Q1 : 1 pt ; Q2 : 1 pt ; Q3 : 1 pt ; Q4 : 1 pt}

\textbf{1.} $315° = 315 \times \dfrac{\pi}{180} = \dfrac{7\pi}{4}$.

\textbf{2.} $\dfrac{5\pi}{3} = \dfrac{5\pi}{3} \times \dfrac{180}{\pi} = 300°$.

\textbf{3.} $-\dfrac{13\pi}{6} + 2\pi = -\dfrac{13\pi}{6} + \dfrac{12\pi}{6} = -\dfrac{\pi}{6} \in {]-\pi;\pi]}$. Mesure principale : $-\dfrac{\pi}{6}$.

\textbf{4.} $\dfrac{3\pi}{4} = \pi - \dfrac{\pi}{4}$ : $\cos\dfrac{3\pi}{4} = -\cos\dfrac{\pi}{4} = -\dfrac{\sqrt{2}}{2}$, $\sin\dfrac{3\pi}{4} = \sin\dfrac{\pi}{4} = \dfrac{\sqrt{2}}{2}$.

Coordonnees : $\left(-\dfrac{\sqrt{2}}{2}\,;\,\dfrac{\sqrt{2}}{2}\right)$.

%===============================================================
\section*{Exercice 2 \hfill (4 points) — C2}
%===============================================================

\baremeIndicatif{Q1 : 1,5 pt ; Q2 : 1 pt ; Q3 : 1,5 pt}

\textbf{1.} $\dfrac{4\pi}{3} = \pi + \dfrac{\pi}{3}$ :
$\cos\dfrac{4\pi}{3} = -\cos\dfrac{\pi}{3} = -\dfrac{1}{2}$, \quad $\sin\dfrac{4\pi}{3} = -\sin\dfrac{\pi}{3} = -\dfrac{\sqrt{3}}{2}$.

\textbf{2.} $\dfrac{11\pi}{6} = 2\pi - \dfrac{\pi}{6}$ :
$\cos\dfrac{11\pi}{6} = \cos\dfrac{\pi}{6} = \dfrac{\sqrt{3}}{2}$.

\textbf{3.} $\cos^2 x = 1 - \dfrac{25}{169} = \dfrac{144}{169}$. Comme $x \in \left]0;\dfrac{\pi}{2}\right[$, $\cos x > 0$, donc $\cos x = \dfrac{12}{13}$.

%===============================================================
\section*{Exercice 3 \hfill (4 points) — C1}
%===============================================================

\baremeIndicatif{Q1 : 2 pts ; Q2 : 1 pt ; Q3 : 1 pt}

\textbf{1.} Avec $s=r\theta$, $s=8\times\dfrac{3\pi}{4}=\boxed{6\pi\ \text{cm}}$.

\textbf{2.} $\theta=\dfrac{s}{r}=\dfrac74=\boxed{\dfrac74\ \text{rad}}$.

\textbf{3.} La fraction de tour vaut $\dfrac{5\pi/3}{2\pi}=\boxed{\dfrac56}$.

%===============================================================
\section*{Exercice 4 \hfill (4 points) — C2}
%===============================================================

\baremeIndicatif{Q1 : 1,5 pt ; Q2 : 1,5 pt ; Q3 : 1 pt}

\textbf{1.} Le point image de $\dfrac{4\pi}{3}$ a pour coordonnees $\boxed{\left(-\dfrac12;-\dfrac{\sqrt3}{2}\right)}$.

\textbf{2.} Le point image de $\dfrac{7\pi}{4}$ a pour coordonnees $\boxed{\left(\dfrac{\sqrt2}{2};-\dfrac{\sqrt2}{2}\right)}$.

\textbf{3.} Oui : la lecture du cercle a l'angle $5\pi/6$ donne exactement les coordonnees de $N$.

%===============================================================
\section*{Exercice 5 \hfill (4 points) — C1, C2}
%===============================================================

\baremeIndicatif{Q1 : 1 pt ; Q2 : 1,5 pt ; Q3 : 1 pt ; Q4 : 0,5 pt}

\textbf{1.} La longueur parcourue vaut $18\times\dfrac{3\pi}{4}=\boxed{\dfrac{27\pi}{2}\ \text{m}}$.

\textbf{2.} Le point associe a $3\pi/4$ a pour coordonnees $\boxed{\left(-\dfrac{\sqrt2}{2};\dfrac{\sqrt2}{2}\right)}$.

\textbf{3.} La hauteur est $\boxed{20+18\times\dfrac{\sqrt2}{2}=20+9\sqrt2\ \text{m}}$.

\textbf{4.} $20+9\sqrt2\approx\boxed{32{,}7\ \text{m}}$.

\end{corrige}
